\section{Proof details for Section~\ref{sec:composition}}\label{app:composition}

\subsection{Lower bounds for Theorem~\ref{thm:interleave}}\label{app:interleave}

Let $\nu_A$ and $\nu_C$ be probability measures on $[\ell,u]$ with equal moments
of orders one and two. For any polynomial $p$ of degree at most two,
$\int\dist(y,A)^2d\nu_A+\int\dist(y,C)^2d\nu_C$ equals
$\int(\dist(y,A)^2-p)\,d\nu_A+\int(\dist(y,C)^2+p)\,d\nu_C$, so it suffices to
find $p$ with $\dist(\cdot,A)^2-p\ge\eta_A$ and $\dist(\cdot,C)^2+p\ge\eta_C$ on
$\R$ and $\eta_A+\eta_C=\delta^2/2$. Because $\dist(y,S)^2=\min_{s\in S}(y-s)^2$,
each inequality can be checked one point $s$ at a time. We use
$\inf_y[\alpha(y-s)^2+\beta(y-t)^2]=\alpha\beta(s-t)^2/(\alpha+\beta)$ for
$\alpha+\beta>0$.

\emph{Separated sets.} By symmetry let $\max A<\min C$, put $a_0=\max A$,
$c_0=\min C$, $\delta=c_0-a_0$, $t_0=(a_0+c_0)/2$ and $p(y)=\delta(y-t_0)$. For
$s\in A$, $\inf_y[(y-s)^2-\delta(y-t_0)]=-\delta^2/4-\delta(s-t_0)\ge\delta^2/4$, and
for $t\in C$, $\inf_y[(y-t)^2+\delta(y-t_0)]=-\delta^2/4+\delta(t-t_0)\ge\delta^2/4$.

\emph{Nested sets.} By symmetry let $C$ lie strictly between the two points
$a_-<a_+$ of $A$. Put $t_0=(\min C+\max C)/2$, $h=(\max C-\min C)/2$ and
$\rho=\delta+h$; then $\abs{s-t_0}\ge\rho$ for $s\in A$ and $\abs{t-t_0}\le h$ for
$t\in C$. If $h>0$, let $\gamma=(\rho-h)/(\rho+h)\in(0,1)$ and
$p(y)=-\gamma(y-t_0)^2$. Then
$\inf_y[(y-s)^2+\gamma(y-t_0)^2]=\gamma(s-t_0)^2/(1+\gamma)\ge\gamma\rho^2/(1+\gamma)$ and
$\inf_y[(y-t)^2-\gamma(y-t_0)^2]=-\gamma(t-t_0)^2/(1-\gamma)\ge-\gamma h^2/(1-\gamma)$.
With $1+\gamma=2\rho/(\rho+h)$ and $1-\gamma=2h/(\rho+h)$ the two bounds add up to
$(\rho-h)^2/2=\delta^2/2$. If $h=0$, take $p(y)=-(y-t_0)^2$; then
$\dist(\cdot,C)^2+p\equiv0$ and $\dist(\cdot,A)^2-p\ge\delta^2/2$.

\subsection{A separating polynomial for Theorem~\ref{thm:alternation}}\label{app:alternation}

Let $A$ and $C$ be disjoint, nonempty and compact, let $T=A\cup C$, and define
the label $\epsilon(t)=+1$ for $t\in A$ and $\epsilon(t)=-1$ for $t\in C$. Let
$\varepsilon=\dist(A,C)>0$. Starting from $s_0=\min T$, let $s_{i+1}$ be the
smallest point of $T$ larger than $s_i$ with the opposite label, as long as
one exists. Such a point exceeds $s_i$ by at least $\varepsilon$, so the
process stops at some $s_r$, and every point of $T$ in $[s_i,s_{i+1})$ has the
label of $s_i$. Every alternating chain visits these runs in increasing
order, so $\operatorname{alt}(A,C)=r$. Let $e_i$ be the largest point of $T$
in $[s_{i-1},s_i)$ and choose $\rho_i\in(e_i,s_i)$. Then
$p(t)=\epsilon(s_0)\prod_{i=1}^r(\rho_i-t)$ has degree $r$, and a point of the
$i$-th run lies between $\rho_i$ and $\rho_{i+1}$, so $\operatorname{sign}p(t)=\epsilon(t)$:
$p>0$ on $A$ and $p<0$ on $C$. If $r\le\kappa$, then $\int p\,d\mu_A=\int p\,d\mu_C$
for any probability measures $\mu_A$ on $A$ and $\mu_C$ on $C$ with equal
moments up to order $\kappa$, which is impossible.

\subsection{Dense moment matrices for the path family}\label{app:sdp}

We use the notation of Theorem~\ref{thm:interleave}, assume that $A$ and $C$
interleave, and let $v\in R$ be a point with $\operatorname{lin}\Phi(v)=0$. Let
$M$ be the moment matrix of $(1,x,y,z)$ with the entries of $v$ and an
unspecified entry $p_{xz}$.

\emph{The dense moment matrix alone does not exclude $v$.} The specified
entries of $M$ form a chordal pattern whose two maximal blocks, the moment
matrices of $(1,x,y)$ and $(1,y,z)$, are positive semidefinite because the two
parts of $v$ are moment vectors. Hence $M$ has a positive semidefinite
completion \citep{GroneEtAl1984}, and $v$ together with this completion
satisfies the dense moment constraint and all RLT inequalities of the edges
and squares.

\emph{The McCormick inequalities of $xz$ alone do not exclude $v$.} They
require $\max\{0,m_x+m_z-1\}\le p_{xz}\le\min\{m_x,m_z\}$, an interval that is
nonempty for $m_x,m_z\in[0,1]$.

\emph{Together they exclude it.} Replacing $x$ by $1-x$ exchanges $a_1$ and $a_2$
and maps the four McCormick inequalities of $xz$ to each other, and similarly
for $z$, so we may assume $a_1<a_2$ and $c_1<c_2$. The $(x,y)$ part of $v$ is
the moment vector of a measure supported on the zero set of $\Phi_A$, which is
$\{(0,a_1),(1,a_2)\}$: at a point with $0<x<1$, concavity or affinity in $x$
gives $\Phi_A(x,y)\ge(1-x)(y-a_1)^2+x(y-a_2)^2>0$. Hence $x=\alpha+\beta y$ on this
support, with $\beta=1/(a_2-a_1)$ and $\alpha=-a_1\beta$, and the vector
$w=(\alpha,-1,\beta,0)$ satisfies $w\T Mw=0$. If $M\succeq0$, then $Mw=0$, and the
last row gives $p_{xz}=\alpha m_z+\beta p_{yz}$; so the completion is unique. In
the same way $z=(y-c_1)/(c_2-c_1)$ on the $(y,z)$ support, which gives
$m_z=(m_y-c_1)/(c_2-c_1)$ and $p_{yz}=(s_y-c_1m_y)/(c_2-c_1)$. Combining,
\[
p_{xz}=\frac{s_y-(a_1+c_1)m_y+a_1c_1}{(a_2-a_1)(c_2-c_1)}.
\]
The numerator is the expectation of $(y-a_1)(y-c_1)$ under any measure with
moments $(m_y,s_y)$. Evaluating it with the $(x,y)$ measure, which puts mass
$m_x$ on $y=a_2$, gives $p_{xz}=m_x(a_2-c_1)/(c_2-c_1)$; evaluating it with the
$(y,z)$ measure gives $p_{xz}=m_z(c_2-a_1)/(a_2-a_1)$. Since the chords meet at
a point interior to both, $0<m_x<1$ and $0<m_z<1$. If $a_1<c_1<a_2<c_2$, then
$(c_2-a_1)/(a_2-a_1)>1$ and $p_{xz}>m_z$. If $c_1<a_1<c_2<a_2$, then
$(a_2-c_1)/(c_2-c_1)>1$ and $p_{xz}>m_x$. Either way a McCormick inequality is
violated. More strongly, the two ingredients prove $\Phi\ge\delta^2/2$
through~\eqref{eq:sdp-identity}: write
$\xi_A^2+\xi_C^2=\tfrac12(\xi_A+\xi_C)^2+\tfrac12(\xi_A-\xi_C)^2$; the function
$g=\tfrac12(\xi_A-\xi_C)^2+\tfrac12(a_2-a_1)^2x(1-x)+\tfrac12(c_2-c_1)^2z(1-z)$ has no
$x^2$ or $z^2$ terms, so it is bilinear in $(x,z)$ and equals its interpolation
$\sum_{ij}g(i,j)\ell_{ij}$ from the four vertices, where
$g(i,j)=\tfrac12(c_{j+1}-a_{i+1})^2$; and $\sum_{ij}\ell_{ij}=1$.

\subsection{Examples for Proposition~\ref{prop:gain}}\label{app:gain-examples}

With $q_1=\Phi_A$ and $q_2=\Phi_C$ for the nested sets $A=\{0,3\}$ and
$C=\{1,2\}$ on $y\in[0,3]$, both pair minima are zero and
$\Delta=\beta^\ast=1/2$; the re-splitting $q_1+\tfrac12(y-\tfrac32)^2$ and
$q_2-\tfrac12(y-\tfrac32)^2$ has pair minima $\tfrac34$ and $-\tfrac14$, so
$\Delta'=0$, in agreement with the zero gap of Theorem~\ref{thm:interleave}. The
polynomials $p$ of Appendix~\ref{app:interleave} are optimal re-splittings of
this kind. Pairwise intersections of the projections in
Proposition~\ref{prop:gain}(i) do not suffice when $k\ge3$: for three leaves of
the form $\Phi_A$ with $A_1=\{0,1\}$, $A_2=\{1,2\}$ and $A_3=\{0,2\}$ on
$y\in[0,2]$, the projections meet pairwise but have no common point, and
$\Delta=2/3$.
